% DERIVED from formal_layer.json — read-only, do not edit.
\begin{definitionv}{synth_3}[Baseline and potential coordinates]
A potential-data atom consists of a baseline category \(X\), binary potential surrogate measurements \(S(0),S(1)\), and binary potential outcomes \(Y(0),Y(1)\). For a nonnegative integer \(d\), its coordinates range over
\[
X\in\{0,\ldots,d-1\},
\qquad
(S(0),S(1),Y(0),Y(1))\in\{0,1\}^{4},
\]
where the baseline domain is empty when \(d=0\). The collection of all such atoms is finite and has decidable equality.
\end{definitionv}

\begin{definitionv}{synth_4}[Full-data atom coordinates]
A full-data atom, parameterized by a nonnegative integer \(d\), consists of the baseline category \(X\), the potential surrogates \(S(0),S(1)\), the potential outcomes \(Y(0),Y(1)\), the treatment indicator \(A\), the realized surrogate \(S\), the realized outcome \(Y\), and the outcome-arrival indicator \(R\). Its coordinates range over
\[
\bigl(X,S(0),S(1),Y(0),Y(1),A,S,Y,R\bigr)
\in \{0,\ldots,d-1\}\times\{0,1\}^{8},
\]
where the baseline set is empty when \(d=0\). The realized coordinates \(S,Y\) are separate from the potential coordinates; every tuple in the displayed Cartesian product is a full-data atom.
\end{definitionv}

\begin{definitionv}{synth_2}[Observed records and space]
The observed-record space \(O\) consists of tuples
\[
O=(X,A,S,R,RY)\in \{0,\ldots,d-1\}\times\{0,1\}^{4},
\]
where \(d\) is a nonnegative integer and the baseline-label set is empty when \(d=0\). The coordinates are the baseline category \(X\), treatment indicator \(A\), realized surrogate \(S\), outcome-arrival indicator \(R\), and observed outcome mark \(RY\). Each of the last four coordinates is binary. The record space is finite and is equipped with the discrete \(\sigma\)-algebra, so every subset is measurable.
\end{definitionv}

\begin{assumptionv}{ass:iid}[Independent observed sampling]
The observed records \(O_1,\ldots,O_n\), with sample size \(n\), are independent and identically distributed under the observed-data law \(P_O\):
\[
O_1,\ldots,O_n\stackrel{\mathrm{iid}}{\sim}P_O.
\]
\end{assumptionv}

\begin{assumptionv}{ass:randomized-independence}[Randomized treatment independence]
The treatment indicator \(A\) is independent of the baseline category \(X\), the potential surrogates \(S(0),S(1)\), and the potential outcomes \(Y(0),Y(1)\):
\[
A\perp \bigl(X,S(0),S(1),Y(0),Y(1)\bigr).
\]
\end{assumptionv}

\begin{assumptionv}{ass:balanced-randomization}[Balanced treatment assignment]
Under the full-data law \(P\), the treatment indicator \(A\) satisfies
\[
P(A=1)=\frac12.
\]
\end{assumptionv}

\begin{assumptionv}{ass:consistency}[Consistency of realized measurements]
The realized surrogate \(S\) and outcome \(Y\) equal the potential measurements selected by the treatment indicator \(A\):
\[
(S,Y)=(S(A),Y(A)).
\]
\end{assumptionv}

\begin{assumptionv}{ass:mar}[Conditionally independent outcome arrival]
The outcome-arrival indicator \(R\) is conditionally independent of the realized outcome \(Y\), given the baseline category \(X\), treatment indicator \(A\), and realized surrogate \(S\):
\[
R\perp Y\mid(X,A,S).
\]
\end{assumptionv}

\begin{definitionv}{synth_1}[Full-data probability law]
The full-data law \(P\) is a probability mass function on full-data atoms of covariate dimension \(d\), where \(d\) is a nonnegative integer. Each atom consists of a potential-data record and separate binary coordinates \(A,S,Y,R\), representing treatment, the realized surrogate, the realized outcome, and outcome arrival, respectively.
\end{definitionv}

\begin{definitionv}{synth_5}[Cellwise arrival probability $\rho_P$ and outcome mean $\mu_P$]
In the setting of \cref{obj:ass:iid,obj:ass:randomized-independence,obj:ass:balanced-randomization,obj:ass:consistency,obj:ass:mar,obj:ass:arrival-floor,obj:def:law-class,obj:def:identified-functional,obj:def:rate,obj:def:point-risk,obj:def:interval-class,obj:def:length-risk,obj:def:mm-estimator,obj:def:mm-interval,obj:def:paired-prior-handle,obj:lem:identification-infrastructure,obj:lem:upper-risk,obj:lem:parametric-floor-infrastructure,obj:lem:prior-reduction,obj:lem:honest-interval-construction,obj:thm:point-frontier,obj:thm:interval-frontier,obj:prop:complete-arrival-reduction,obj:thm:normalized-converse,obj:synth_1,obj:synth_2,obj:synth_3,obj:synth_4,obj:synth_6}, define the cellwise arrival probability and arrived-outcome conditional mean by
\[
\rho_P(x,a,s)=
\begin{cases}
P(R=1\mid X=x,A=a,S=s),
& P(X=x,A=a,S=s)>0,\\
0,
& P(X=x,A=a,S=s)=0,
\end{cases}
\]
and
\[
\mu_P(x,a,s)=
\begin{cases}
\mathbb E_P[Y\mid X=x,A=a,S=s,R=1],
& P(X=x,A=a,S=s,R=1)>0,\\
0,
& P(X=x,A=a,S=s,R=1)=0.
\end{cases}
\]
Here \(P\) is a full-data law, \(X\in\{0,\ldots,d-1\}\) is the baseline cell, \(A\in\{0,1\}\) is the treatment assignment, \(S\in\{0,1\}\) is the observed surrogate, \(R\in\{0,1\}\) is the outcome-arrival indicator, and \(Y\in\{0,1\}\) is the realized primary outcome. The definitions apply to every \((x,a,s)\in\{0,\ldots,d-1\}\times\{0,1\}^2\).
\end{definitionv}

\begin{assumptionv}{ass:arrival-floor}[Cellwise outcome-arrival floor]
For every cell \((x,a,s)\) with positive probability under the full-data law \(P\), the cellwise arrival probability \(\rho_P(x,a,s)\) is at least the arrival floor \(q\):
\[
P(X=x,A=a,S=s)>0
\quad\Longrightarrow\quad
\rho_P(x,a,s)\ge q.
\]
\end{assumptionv}

\begin{definitionv}{def:law-class}[Law class \(\mathcal M(d,q)\)]
\[
\mathcal M(d,q)
=
\left\{
P:
|\operatorname{supp}_P(X)|\le d,\;
P\text{ satisfies }
\cref{obj:ass:randomized-independence,obj:ass:balanced-randomization,obj:ass:consistency,obj:ass:mar,obj:ass:arrival-floor}
\right\}.
\]
Here \(d\) is a positive integer, \(q\in[1/2,1]\) is the known cellwise arrival floor, and \(\operatorname{supp}_P(X)\) is the set of baseline labels with positive probability under \(P\).
\end{definitionv}

\begin{definitionv}{def:identified-functional}[Observed-data identifying functional]
The observed-data functional is
\[
\Psi(P_O)
=
\sum_{x=0}^{d-1}\sum_{s=0}^{1}
\left\{
\frac{P_O(X=x,A=1,S=s)}{P_O(A=1)}\mu_P(x,1,s)
-
\frac{P_O(X=x,A=0,S=s)}{P_O(A=0)}\mu_P(x,0,s)
\right\},
\]
for any nonnegative integer \(d\) and any observed-record probability law \(P_O\) on
\[
O=\{0,\ldots,d-1\}\times\{0,1\}^{4},
\]
with coordinates \((X,A,S,R,RY)\). The baseline-label set and the outer sum are empty when \(d=0\). The arrived-outcome mean is defined by
\[
\mu_P(x,a,s)
=
\frac{P_O(X=x,A=a,S=s,R=1,RY=1)}
     {P_O(X=x,A=a,S=s,R=1)},
\qquad a,s\in\{0,1\}.
\]
Every ratio with zero denominator is assigned the value zero. Thus each weighted contribution is zero whenever \(P_O(X=x,A=a,S=s)=0\).
\end{definitionv}

\begin{definitionv}{synth_6}[Measurable estimator class $\mathcal T_n$]
For integers \(n,d\ge1\), define the measurable estimator class by
\[
\mathcal T_n
=
\left\{
T:\bigl(\{0,\ldots,d-1\}\times\{0,1\}^4\bigr)^n
\longrightarrow[-1,1]
:\ T\text{ is measurable}
\right\}.
\]
Here the observed record \(O=(X,A,S,R,RY)\) takes values in the Cartesian observed-record space \(\{0,\ldots,d-1\}\times\{0,1\}^4\), with the baseline-label convention of \cref{obj:synth_2}. Thus each \(T\in\mathcal T_n\) is defined on every sample in this Cartesian space and takes values in the treatment-effect range \([-1,1]\), in the setting of \cref{obj:ass:iid,obj:ass:randomized-independence,obj:ass:balanced-randomization,obj:ass:consistency,obj:ass:mar,obj:ass:arrival-floor,obj:def:law-class,obj:def:identified-functional,obj:def:rate,obj:def:point-risk,obj:def:interval-class,obj:def:length-risk,obj:def:mm-estimator,obj:def:mm-interval,obj:def:paired-prior-handle,obj:lem:identification-infrastructure,obj:lem:upper-risk,obj:lem:parametric-floor-infrastructure,obj:lem:prior-reduction,obj:lem:honest-interval-construction,obj:thm:point-frontier,obj:thm:interval-frontier,obj:prop:complete-arrival-reduction,obj:thm:normalized-converse,obj:synth_1,obj:synth_3,obj:synth_4,obj:synth_5}.
\end{definitionv}

\begin{definitionv}{def:point-risk}[Minimax squared-error risk \(\mathfrak R^*_{n,d,q}\)]
\[
\mathfrak R^*_{n,d,q}
=
\inf_{T\in\mathcal T_n}
\sup_{P\in\mathcal M(d,q)}
\mathbb E_P
\left[
\left\{
T(O_1,\ldots,O_n)-\tau(P)
\right\}^2
\right].
\]
\end{definitionv}

\begin{definitionv}{def:interval-class}[Honest connected confidence intervals]
The class \(\mathcal C_\alpha(n,d,q)\), for nonnegative integers \(n,d\) and real parameters \(q,\alpha\), is defined using the observed-record space \(O\) with the baseline-label convention of \cref{obj:synth_2}. It consists of procedures \(I(o)=[L(o),U(o)]\) with measurable endpoints satisfying
\[
-1\le L(o)\le U(o)\le1
\qquad\text{for every }o\in O^n,
\]
and uniform coverage:
\[
\mathcal C_\alpha(n,d,q)
=
\left\{
I:
P_O^{\otimes n}\!\left\{
o\in O^n:\tau(P)\in[L(o),U(o)]
\right\}\ge1-\alpha
\quad\text{for every }P\in\mathcal M(d,q)
\right\}.
\]
Here \(\mathcal M(d,q)\) is the law class in \cref{obj:def:law-class}, and \(\tau(P)=\mathbb E_P[Y(1)-Y(0)]\) is the superpopulation average treatment effect. For \(q\in[1/2,1]\) and \(P\in\mathcal M(d,q)\), \cref{obj:lem:identification-infrastructure} gives \(\tau(P)=\Psi(P_O)\), with \(\Psi\) defined in \cref{obj:def:identified-functional}.
\end{definitionv}

\begin{definitionv}{def:length-risk}[Minimax expected length \(\mathfrak L^*_{n,d,q,\alpha}\)]
\[
\mathfrak L^*_{n,d,q,\alpha}
=
\inf_{I\in\mathcal C_\alpha(n,d,q)}
\sup_{P\in\mathcal M(d,q)}
\mathbb E_P
\left[
\operatorname{length}\{I(O^n)\}
\right].
\]
For \(I(O^n)=[L,U]\), its length is \(U-L\).
\end{definitionv}

\begin{definitionv}{def:rate}[Squared-error benchmark \(r_{n,d,q}\)]
\[
r_{n,d,q}
=
\frac{1}{n}
+
(1-q)^2
\min\left\{
1,
\left[\frac{d}{n\log(e+n)}\right]^2
\right\}
=
\frac{1}{n}+g_{n,d,q}^2.
\]
\end{definitionv}

\begin{algorithmv}{def:mm-estimator}[Moment estimator $\widehat\tau_{\mathrm{MM}}$]
Given observed records \(O_1,\ldots,O_n\), sample size \(n\), support bound \(d\), and arrival floor \(q\), define \(\widehat\tau_{\mathrm{MM}}\) using the tuning parameters \(m,k,B\) as follows.
\begin{enumerate}
\item Set
\[
\ell_n=\log(e+n),\qquad
m=\frac n8,\qquad
k=\max\{1,\lfloor\ell_n/64\rfloor\},\qquad
B=256\ell_n.
\]
Write the supplied fifth coordinate of each observed record as
\[
Y_i^{\mathrm{obs}}=(RY)_i.
\]
On the support of the observation map, this coordinate equals \(R_iY_i\) and is zero when the outcome is missing.

\item Draw
\[
N\sim\operatorname{Poisson}(n/2).
\]
If \(N>n\), set \(\widetilde\tau=0\). If \(N\le n\), independently assign each of the first \(N\) records a uniform stream indicator and define the stream index sets by
\[
Z_i\in\{1,2,3,4\},\qquad
\Pr(Z_i=t)=\frac14,\qquad
\mathcal I_t=\{i\in\{1,\ldots,N\}:Z_i=t\}.
\]
For this branch, perform the following three steps.

\item For each cell \(j=(x,a,s)\), write \(a_j=a\) and define
\[
M_j^{(2)}
=\sum_{i\in\mathcal I_2}
\mathbf 1\{(X_i,A_i,S_i)=j,\ R_i=0\},
\qquad
V_j=\frac{M_j^{(2)}}m,
\]
\[
C'_j
=\sum_{i\in\mathcal I_3}
\mathbf 1\{(X_i,A_i,S_i)=j,\ R_i=1\},
\qquad
C_j
=\sum_{i\in\mathcal I_4}
\mathbf 1\{(X_i,A_i,S_i)=j,\ R_i=1\},
\]
\[
U_j
=\sum_{i\in\mathcal I_4}
\mathbf 1\{(X_i,A_i,S_i)=j,\ Y_i^{\mathrm{obs}}=1\}.
\]
Thus \(U_j\) is the arrived-success count on the support of the observation map.

\item Let \(T_k\) be the degree-\(k\) Chebyshev polynomial, with convention
\[
T_k(\cos\theta)=\cos(k\theta).
\]
Define
\[
Q_k(z)=
\begin{cases}
\dfrac{1-T_k(1-2z/B)}{2k^2z/B},&z\ne0,\\[6pt]
1,&z=0,
\end{cases}
\qquad
1-Q_k(z)=\sum_{v=1}^{k-1}a_vz^v.
\]
Using the falling-factorial convention
\[
(C_j-1)_{v-1}=\prod_{t=1}^{v-1}(C_j-t),
\]
where an empty product equals one, set
\[
H_j=\sum_{v=1}^{k-1}a_v(U_j-C_j/2)(C_j-1)_{v-1},
\qquad
D_j=
\begin{cases}
U_j/C_j-1/2,&C_j>0,\\
0,&C_j=0,
\end{cases}
\]
\[
G_j
=H_j\mathbf 1\{C'_j\le B/4\}
+D_j\mathbf 1\{C'_j>B/4\}.
\]

\item Define the interval projection and randomized estimate by
\[
\Pi_{[-1,1]}(t)=\max\{-1,\min\{1,t\}\},
\]
\[
\widetilde\tau
=\Pi_{[-1,1]}\!\left[
\frac2m\sum_{i\in\mathcal I_1}(2A_i-1)R_i
\left(Y_i^{\mathrm{obs}}-\frac12\right)
+2\sum_j(2a_j-1)V_jG_j
\right].
\]

\item Output
\[
\widehat\tau_{\mathrm{MM}}
=
\begin{cases}
\displaystyle
\Pi_{[-1,1]}\!\left[
\frac2n\sum_{i=1}^n(2A_i-1)Y_i^{\mathrm{obs}}
\right],
&q=1,\\[8pt]
\displaystyle
\Pi_{[-1,1]}\!\left[
\frac2n\sum_{i=1}^n(2A_i-1)R_i
\left(Y_i^{\mathrm{obs}}-\frac12\right)
\right],
&q\ne1\ \text{and}\ (\ell_n<128\ \text{or}\ d\ge n\ell_n),\\[8pt]
\displaystyle
\mathbb E[\widetilde\tau\mid O_1,\ldots,O_n],
&q\ne1,\ \ell_n\ge128,\ d<n\ell_n.
\end{cases}
\]
The conditional expectation averages over the auxiliary Poisson draw and stream assignments.
\end{enumerate}
\end{algorithmv}

\begin{definitionv}{def:mm-interval}[Confidence interval $\widehat I_{\mathrm{MM},\alpha}$]
Define
\[
\widehat I_{\mathrm{MM},\alpha}(O^n)
=
\left[
\max\{-1,\widehat\tau_{\mathrm{MM}}-h_{n,d,q,\alpha}\},
\min\{1,\widehat\tau_{\mathrm{MM}}+h_{n,d,q,\alpha}\}
\right],
\]
where the clipped confidence radius is
\[
h_{n,d,q,\alpha}
=\min\left\{2,\sqrt{\frac{C_Ur_{n,d,q}}{\alpha}}\right\},
\]
and \(C_U\) is any universal constant for which the upper squared-risk bound for \(\widehat\tau_{\mathrm{MM}}\) holds.
\end{definitionv}

\begin{theoremv}{thm:point-frontier}[Finite-sample minimax risk frontier]
There exist universal numerical constants \(0<c<C<\infty\) such that, for every pair of positive integers \(n,d\) and every known arrival floor \(q\in[1/2,1]\),
\[
c\,r_{n,d,q}\le \mathfrak R^*_{n,d,q}\le C\,r_{n,d,q},
\]
where \(r_{n,d,q}\) is the squared-error benchmark defined in \cref{obj:def:rate} and \(\mathfrak R^*_{n,d,q}\) is the minimax squared-error risk defined in \cref{obj:def:point-risk}.
\end{theoremv}

\begin{theoremv}{thm:interval-frontier}[Fixed coverage length frontier]
For each fixed noncoverage level \(\alpha\in(0,1/2)\), there exist constants \(0<c_\alpha<C_\alpha<\infty\), depending only on \(\alpha\), such that, for every pair of positive integers \(n,d\) and every \(q\in[1/2,1]\), the minimax expected interval length \(\mathfrak L^*_{n,d,q,\alpha}\) from \cref{obj:def:length-risk} and the squared-error benchmark \(r_{n,d,q}\) from \cref{obj:def:rate} satisfy
\[
c_\alpha\sqrt{r_{n,d,q}}
\le \mathfrak L^*_{n,d,q,\alpha}
\le C_\alpha\sqrt{r_{n,d,q}}.
\]
\end{theoremv}

\begin{propositionv}{prop:complete-arrival-reduction}[Complete arrival and parametric frontiers]
There exist universal constants \(0<c<C\) such that, for every noncoverage level \(\alpha\in(0,1/2)\), there exist constants \(0<c_\alpha<C_\alpha\), depending only on \(\alpha\), for which the following statements hold for every positive integer sample size \(n\) and positive integer baseline-support bound \(d\) at arrival floor \(q=1\).

The error scale and squared-error benchmark satisfy
\[
g_{n,d,1}
=(1-1)\min\left\{1,\frac{d}{n\log(e+n)}\right\}=0,
\qquad
r_{n,d,1}=\frac1n,
\]
where the benchmark is defined in \cref{obj:def:rate}.
For every sample \(O^n=(O_1,\ldots,O_n)\) in the Cartesian observed-record space, the estimator of \cref{obj:def:mm-estimator} satisfies
\[
\widehat\tau_{\mathrm{MM}}(O^n)
=
\Pi_{[-1,1]}
\left[
\frac{2}{n}\sum_{i=1}^{n}(2A_i-1)Y_i^{\mathrm{obs}}
\right],
\]
where \(A_i\) is the binary treatment coordinate, \(Y_i^{\mathrm{obs}}=(RY)_i\) is the supplied fifth coordinate of record \(i\), and
\(\Pi_{[-1,1]}(z)=\max\{-1,\min\{1,z\}\}\) is interval projection.

For every full-data law \(P\in\mathcal M(d,1)\) and every such sample, with population average treatment effect
\(\tau(P)=\mathbb E_P[Y(1)-Y(0)]\), where \(Y(1)\) and \(Y(0)\) are the binary potential outcomes,
\[
\left|\widehat\tau_{\mathrm{MM}}(O^n)-\tau(P)\right|
\le
\left|
\frac{2}{n}\sum_{i=1}^{n}(2A_i-1)Y_i^{\mathrm{obs}}-\tau(P)
\right|.
\]
Under the iid observed-sample law induced by \(P\),
\[
\sup_{P\in\mathcal M(d,1)}
\mathbb E_P\!\left[
\bigl(\widehat\tau_{\mathrm{MM}}(O^n)-\tau(P)\bigr)^2
\right]
\le
\sup_{P\in\mathcal M(d,1)}
\mathbb E_P\!\left[
\left(
\frac{2}{n}\sum_{i=1}^{n}(2A_i-1)Y_i^{\mathrm{obs}}-\tau(P)
\right)^2
\right]
\le \frac4n.
\]
The minimax squared-error risk and minimax worst-law expected honest interval length of
\cref{obj:def:point-risk,obj:def:length-risk} satisfy, uniformly over \(n,d\ge1\),
\[
\frac{c}{n}
\le \mathfrak R^*_{n,d,1}
\le \frac{C}{n},
\qquad
\frac{c_\alpha}{\sqrt n}
\le \mathfrak L^*_{n,d,1,\alpha}
\le \frac{C_\alpha}{\sqrt n}.
\]
\end{propositionv}

\begin{lemmav}{lem:identification-infrastructure}[Randomized MAR identification]
Let \(d\) be a nonnegative integer, let \(q\in[1/2,1]\) be the cellwise arrival floor, and let \(P\) be a full-data probability law with baseline variable \(X\) taking values in a set of \(d\) labels, binary treatment \(A\), binary potential surrogates \(S(0),S(1)\), binary potential outcomes \(Y(0),Y(1)\), binary realized surrogate and outcome \(S,Y\), and binary outcome-arrival indicator \(R\). Suppose that \(P\) satisfies:
\begin{itemize}
\item \textbf{(Randomized independence.)} \cref{obj:ass:randomized-independence}.
\item \textbf{(Balanced randomization.)} \cref{obj:ass:balanced-randomization}.
\item \textbf{(Consistency.)} \cref{obj:ass:consistency}.
\item \textbf{(Missing at random.)} \cref{obj:ass:mar}.
\item \textbf{(Arrival floor.)} \cref{obj:ass:arrival-floor} with floor \(q\).
\end{itemize}
Write \(P_O\) for the image of \(P\) under the observation map \(O=(X,A,S,R,RY)\), where \(RY\) is the observed outcome mark. Then the superpopulation average treatment effect satisfies
\[
\tau(P)=\mathbb E_P[Y(1)-Y(0)]=\Psi(P_O),
\]
where \(\Psi\) is the observed-law functional in \cref{obj:def:identified-functional}, with every contribution from a null \((X,A,S)\) cell equal to zero by definition.
\end{lemmav}

\begin{lemmav}{lem:upper-risk}[Uniform upper risk bound]
There exists a universal numerical constant \(0<C_U<\infty\) such that, for every choice satisfying
\begin{itemize}
\item \textbf{(Sample size and dimension.)} \(n,d\) are integers with \(n\ge1\) and \(d\ge1\);
\item \textbf{(Arrival floor.)} \(q\in[1/2,1]\);
\item \textbf{(Law and sampling.)} \(P\in\mathcal M(d,q)\), and \(O_1,\ldots,O_n\) are independent observed records with common law \(P_O\), the observed-law image of \(P\);
\end{itemize}
the estimator \(\widehat\tau_{\mathrm{MM}}\) satisfies
\[
\mathbb E_P\!\left[
  \bigl(\widehat\tau_{\mathrm{MM}}-\tau(P)\bigr)^2
\right]
\le C_U r_{n,d,q},
\]
where the expectation is over the stated independent sample.
\end{lemmav}

\begin{lemmav}{lem:honest-interval-construction}[Honest clipped confidence intervals]
Suppose that:
\begin{itemize}
\item \textbf{(Sample size and dimension.)} \(n,d\) are integers with \(n,d\ge 1\).
\item \textbf{(Arrival floor.)} \(q\in[1/2,1]\).
\item \textbf{(Noncoverage level.)} \(\alpha\in(0,1/2)\).
\item \textbf{(Uniform risk certificate.)} \(C_U>0\) certifies the uniform squared-error bound
\[
\mathbb E_{P_O^{\otimes n}}
\left[
  \bigl(\widehat\tau_{\mathrm{MM}}(O^n)-\tau(P)\bigr)^2
\right]
\le C_U r_{n,d,q}
\]
for every \(n,d\ge1\), every \(q\in[1/2,1]\), and every \(P\in\mathcal M(d,q)\), under any iid observed-sample law. Here \(r_{n,d,q}\) is defined in \cref{obj:def:rate}, and
\(\tau(P)=\mathbb E_P[Y(1)-Y(0)]\) is the population average treatment effect.
\end{itemize}
Let \(\widehat I_{\mathrm{MM},\alpha}\) and its radius \(h_{n,d,q,\alpha}\) be as in \cref{obj:def:mm-interval}, using this \(C_U\). The interval has uniform coverage:
\[
P_O^{\otimes n}
\left\{
  \tau(P)\in\widehat I_{\mathrm{MM},\alpha}(O^n)
\right\}
\ge 1-\alpha
\qquad
\text{for every }P\in\mathcal M(d,q).
\]
For every observed sample \(O^n\), its length satisfies
\[
\operatorname{length}
\left\{\widehat I_{\mathrm{MM},\alpha}(O^n)\right\}
\le
\min\{2,2h_{n,d,q,\alpha}\}
=
\min\left\{2,2\sqrt{\frac{C_Ur_{n,d,q}}{\alpha}}\right\}.
\]
Consequently, for every \(P\in\mathcal M(d,q)\),
\[
\mathbb E_{P_O^{\otimes n}}
\left[
  \operatorname{length}
  \left\{\widehat I_{\mathrm{MM},\alpha}(O^n)\right\}
\right]
\le
\min\left\{2,2\sqrt{\frac{C_Ur_{n,d,q}}{\alpha}}\right\}.
\]
If \(h_{n,d,q,\alpha}<2\), then for every \(P\in\mathcal M(d,q)\),
\[
P_O^{\otimes n}
\left\{
  \tau(P)\notin\widehat I_{\mathrm{MM},\alpha}(O^n)
\right\}
\le\alpha.
\]
If \(h_{n,d,q,\alpha}=2\), then for every \(P\in\mathcal M(d,q)\) and every observed sample \(O^n\),
\[
\tau(P)\in\widehat I_{\mathrm{MM},\alpha}(O^n).
\]
\end{lemmav}

\begin{definitionv}{def:paired-prior-handle}[Normalized paired prior construction]
The paired-prior handle consists of the normalized priors \(\Pi_-,\Pi_+\), their observed-sample mixtures \(Q_-^{(n)},Q_+^{(n)}\), their associated unnormalized count laws, a fixed fallback record \(o_{\circ}\), and the count-ordering kernel \(\kappa_n\), constructed below for parameters satisfying:
\begin{itemize}
\item \textbf{(Sample size.)} \(n\) is an integer with \(n\ge1\).
\item \textbf{(Baseline dimension.)} \(d\) is an integer with \(d\ge1\).
\item \textbf{(Arrival floor.)} \(q\in[1/2,1]\).
\end{itemize}
The observed-record space \(O\) consists of tuples \((X,A,S,R,RY)\), with \(X\) taking values in a fixed set of \(d\) baseline labels and the remaining coordinates binary. Fix \(o_{\circ}\in O\).

Set
\[
\delta=1-q,\qquad
L=\log(e+n),\qquad
K=\lceil8L\rceil,\qquad
h=K^{-2},\qquad
m=K-1,\qquad
B_0=\frac{K}{1024n},
\]
\[
M=\min\left\{
\left\lfloor\frac{d-1}{2}\right\rfloor,
\lfloor nL\rfloor
\right\},
\qquad
c_0=\frac1{16},\qquad
b=1-\frac{\delta}{2},\qquad
\mu_0=\frac{b^2}{2q}.
\]
For \(0\le i\le m\), define the interpolation nodes and their Lagrange cardinal polynomials by
\[
x_i=\frac{1+h}{2}+\frac{1-h}{2}\cos\frac{i\pi}{m},
\qquad
l_i(t)=\prod_{\substack{0\le t'\le m\\t'\ne i}}
\frac{t-x_{t'}}{x_i-x_{t'}}.
\]
Set
\[
D=\sum_{i=0}^{m}|l_i(0)|,
\qquad
\nu_i=\frac{l_i(0)}{D}.
\]
The latent pair \((p,z)\in[0,B_0]\times\{-1,1\}\) has distribution
\[
\Pr\{(p,z)=(B_0x_i,\operatorname{sign}(\nu_i))\}
=\frac{h|\nu_i|}{x_i},
\qquad 0\le i\le m,
\]
\[
\Pr\{(p,z)=(0,1)\}
=1-\sum_{i=0}^{m}\frac{h|\nu_i|}{x_i}.
\]
Draw independent copies \((p_j,z_j)\), \(1\le j\le M\), and independent fair orientation signs, independent also of these copies, with
\[
\Pr(r_j=1)=\Pr(r_j=-1)=\frac12.
\]

For each \(\sigma\in\{-1,1\}\), pair \(j\) consists of two distinct baseline labels, each with unnormalized mass \(p_j\), with orientations
\[
u\in\{r_j,-r_j\}.
\]
At each label, \(X\) equals that label, and the potential surrogates and control outcome satisfy
\[
S(0)=S(1)=0,\qquad Y(0)=0.
\]
Treatment \(A\) is an independent fair coin, and control arrival is complete. The treated arrival probability and treated outcome mean are
\[
\rho_{\sigma,j,u}
=1-\frac{\delta}{2}(1+\sigma z_ju),
\qquad
\mu_{\sigma,j,u}
=\frac{b/2+c_0u}{\rho_{\sigma,j,u}}.
\]
Draw \(Y(1)\) as a Bernoulli variable with mean \(\mu_{\sigma,j,u}\), and draw treated arrival independently of \(Y(1)\) with probability \(\rho_{\sigma,j,u}\). In the coordinate order
\[
(A,R,RY)=(0,1,0),(1,0,0),(1,1,1),(1,1,0),
\]
the four observed atom coefficients at orientation \(u\) are
\[
v_{\sigma}(z_j,u)
=
\left(
\frac12,
\frac{\delta(1+\sigma z_ju)}4,
\frac b4+\frac{c_0u}{2},
\frac b4-\frac{c_0u}{2}-\frac{\sigma\delta z_ju}{4}
\right).
\]
The two orientations specify the complete eight-coordinate observed atom vector for pair \(j\).

Add one distinct fixed filler label with unnormalized mass
\[
f=1-2MhB_0.
\]
At this label, set
\[
S(0)=S(1)=0,\qquad Y(0)=0,
\]
let treatment be an independent fair coin, make arrival complete in both treatment arms, and draw \(Y(1)\) as a Bernoulli variable with mean \(\mu_0\). Its observed atom coefficients, in the same coordinate order, are
\[
v_{\mathrm{fill}}
=\left(\frac12,0,\frac{\mu_0}{2},\frac{1-\mu_0}{2}\right).
\]
At every label, the realized variables and observed outcome mark satisfy
\[
S=S(A),\qquad Y=Y(A),\qquad RY=R\,Y.
\]

Normalize the label masses by
\[
J=f+2\sum_{j=1}^{M}p_j,
\qquad
f\longmapsto\frac fJ,
\qquad
p_j\longmapsto\frac{p_j}{J}
\]
for each of the two labels in pair \(j\). Let \(P_{\sigma}\) be the resulting normalized causal law and \(P_{\sigma,O}\) its observed-law image. Define the prior and complete observed-sample mixture by
\[
\Pi_{\sigma}=\operatorname{Law}(P_{\sigma}),
\qquad
Q_{\sigma}^{(n)}
=\int P_{\sigma,O}^{\otimes n}\,d\Pi_{\sigma}(P_{\sigma}).
\]

For the associated unnormalized count experiment, set
\[
\lambda=4n.
\]
Conditional on the latent variables, assign independent Poisson counts to the eight atoms of each pair and the filler atoms, with respective mean vectors
\[
\lambda p_jv_{\sigma}(z_j,r_j),
\qquad
\lambda p_jv_{\sigma}(z_j,-r_j),
\qquad
\lambda f v_{\mathrm{fill}}.
\]
All other observed atoms have count zero. The associated count law is the mixture of this conditional count experiment over the latent draws.

Define the Markov kernel \(\kappa_n\) by uniformly ordering the counted multiset and returning its first \(n\) records when the total count is at least \(n\); otherwise it returns
\[
(o_{\circ},\ldots,o_{\circ})\in O^n.
\]
The kernel \(\kappa_n\) is independent of \(\sigma\) and of every latent variable.
\end{definitionv}

\begin{theoremv}{thm:normalized-converse}[Normalized priors with separated effects]
For every fixed \(\beta\in(0,1/4)\), there exists a constant \(c_\beta>0\) such that the following holds for every sample size \(n\), baseline-support bound \(d\), and arrival floor \(q\) satisfying:
\begin{itemize}
\item \textbf{(Sample size and dimension.)} \(n,d\) are integers with \(n,d\ge1\).
\item \textbf{(Arrival floor.)} \(q\in[1/2,1]\).
\item \textbf{(Separation scale.)} The scale
\[
g_{n,d,q}=(1-q)\min\left\{1,\frac{d}{n\log(e+n)}\right\}
\]
satisfies \(g_{n,d,q}^{\,2}\ge n^{-1}\).
\end{itemize}
There exist discrete probability priors \(\Pi_-\) and \(\Pi_+\) supported on \(\mathcal M(d,q)\) from \cref{obj:def:law-class}, whose laws satisfy \cref{obj:ass:randomized-independence,obj:ass:balanced-randomization,obj:ass:consistency,obj:ass:mar,obj:ass:arrival-floor}. Their complete observed-sample mixtures, defined by
\[
Q_-^{(n)}=\int P_O^{\otimes n}\,\Pi_-(dP),
\qquad
Q_+^{(n)}=\int P_O^{\otimes n}\,\Pi_+(dP),
\]
where \(P_O\) is the observed-law image of the full-data law \(P\), satisfy
\[
\operatorname{TV}\!\left(Q_-^{(n)},Q_+^{(n)}\right)
=\sup_{B\subseteq O^n}
\left|Q_-^{(n)}(B)-Q_+^{(n)}(B)\right|
\le\beta.
\]
Here \(O^n\) is the complete observed-record sample space, and the supremum is over measurable events. There also exists a deterministic center \(m_0\in\mathbb R\) such that, for the population average treatment effect
\[
\tau(P)=\mathbb E_P\!\left[Y(1)-Y(0)\right],
\]
where \(Y(1)\) and \(Y(0)\) are the binary potential outcomes,
\[
\Pi_-\!\left\{P:\tau(P)\le m_0-c_\beta g_{n,d,q}\right\}\ge1-\beta,
\qquad
\Pi_+\!\left\{P:m_0+c_\beta g_{n,d,q}\le\tau(P)\right\}\ge1-\beta.
\]
\end{theoremv}

\begin{lemmav}{lem:parametric-floor-infrastructure}[Parametric risk and length floors]
There exists a universal constant \(c_0>0\) such that, for every pair of positive integers \(n,d\) and every \(q\in[1/2,1]\), the minimax squared-error risk \(\mathfrak R^*_{n,d,q}\) defined in \cref{obj:def:point-risk} satisfies
\[
\mathfrak R^*_{n,d,q}\ge \frac{c_0}{n}.
\]
For each noncoverage level \(\alpha\in(0,1/2)\), there exists a constant \(c_{0,\alpha}>0\), depending only on \(\alpha\), such that, for every pair of positive integers \(n,d\) and every \(q\in[1/2,1]\), the minimax worst-law expected honest interval length \(\mathfrak L^*_{n,d,q,\alpha}\) defined in \cref{obj:def:length-risk} satisfies
\[
\mathfrak L^*_{n,d,q,\alpha}\ge \frac{c_{0,\alpha}}{\sqrt n}.
\]
\end{lemmav}

\begin{lemmav}{lem:prior-reduction}[Prior reduction and lower bounds]
Let \(c_\beta\) denote the fixed separation constant selected from the normalized converse at mixture tolerance \(\beta\in(0,1/4)\), and set \(\beta_0=1/16\). Write
\[
\ell_n=\log(e+n),
\qquad
g_{n,d,q}=(1-q)\min\left\{1,\frac{d}{n\ell_n}\right\}.
\]
Use the minimax risks and rate in \cref{obj:def:point-risk,obj:def:length-risk,obj:def:rate}.

There exist universal constants \(c>0\) and \(c_0>0\) such that, for every \(\alpha\in(0,1/2)\), there exist constants \(c_\alpha>0\) and \(c_{0,\alpha}>0\) with the following guarantees. Set
\[
\beta_\alpha=\frac{1-2\alpha}{8}.
\]
For every \(n,d,q\) satisfying
\begin{itemize}
\item \textbf{(Sample size.)} \(n\) is an integer with \(n\ge1\).
\item \textbf{(Support bound.)} \(d\) is an integer with \(d\ge1\).
\item \textbf{(Arrival floor.)} \(q\in[1/2,1]\).
\end{itemize}
whenever \(g_{n,d,q}^2\ge n^{-1}\), both bounds
\[
\mathfrak R^*_{n,d,q}
\ge \frac{13}{32}c_{\beta_0}^{\,2}g_{n,d,q}^{\,2},
\qquad
\mathfrak L^*_{n,d,q,\alpha}
\ge \frac54 c_{\beta_\alpha}(1-2\alpha)g_{n,d,q}
\]
hold. Throughout the stated parameter range, the parametric bounds are
\[
\mathfrak R^*_{n,d,q}\ge\frac{c_0}{n},
\qquad
\mathfrak L^*_{n,d,q,\alpha}\ge\frac{c_{0,\alpha}}{\sqrt n}.
\]
The combined guarantees are
\(\mathfrak R^*_{n,d,q}\ge c\,r_{n,d,q}\) and
\(\mathfrak L^*_{n,d,q,\alpha}\ge c_\alpha\sqrt{r_{n,d,q}}\).
\end{lemmav}
